%% Caractéristique de transfert d'un CAN 3 bits : 8 paliers de largeur q.
%% Ici l'amplitude d'entrée vaut 8 V, donc q = 8 / 2^3 = 1 V.
\documentclass[tikz,border=4pt]{standalone}
\usepackage{liaisons}
\begin{document}
\begin{tikzpicture}[x=0.55cm,y=0.55cm,
  axe/.style={-{Stealth[length=5pt]}, line width=0.6pt},
  grad/.style={font=\scriptsize},
  guide/.style={black!55, dash pattern=on 2pt off 2pt, line width=0.6pt}]

%% ================== repère ============================================
\draw[axe] (-0.3,0) -- (9.1,0);
\node[grad, anchor=north west, inner sep=1pt] at (8.5,-0.55) {$V_e$ (V)};
\draw[axe] (0,-0.3) -- (0,8.3);
\node[grad, anchor=west, inner sep=1pt] at (-1.15,8.75) {Sortie numérique};
\foreach \k/\mot in {0/000,1/001,2/010,3/011,4/100,5/101,6/110,7/111} {
  \draw[line width=0.5pt] (0,\k) -- (-0.18,\k);
  \node[grad, anchor=east, inner sep=2pt] at (-0.2,\k) {\mot}; }
\foreach \k in {1,...,8} {
  \draw[line width=0.5pt] (\k,0) -- (\k,-0.18);
  \node[grad, anchor=north, inner sep=2pt] at (\k,-0.2) {\k}; }

%% ================== droite idéale =====================================
\draw[black!60, line width=0.6pt, dash pattern=on 2.5pt off 2pt] (0.5,0) -- (7.9,7.4);

%% ================== escalier de quantification ========================
\draw[solideB, line width=1.4pt, line join=round]
  (0,0) -- (1,0) -- (1,1) -- (2,1) -- (2,2) -- (3,2) -- (3,3) -- (4,3) --
  (4,4) -- (5,4) -- (5,5) -- (6,5) -- (6,6) -- (7,6) -- (7,7) -- (8,7);

%% ================== cote du quantum q =================================
\draw[guide] (4,0) -- (4,4);
\draw[guide] (5,0) -- (5,4);
\draw[{Stealth[length=4pt]}-{Stealth[length=4pt]}, line width=1pt, solideA] (4,3.5) -- (5,3.5);
\node[grad, text=solideA, anchor=west, inner sep=2pt] at (5.1,3.5) {$q$};

%% ================== rappel de la formule ==============================
\node[draw=black!45, line width=0.6pt, rounded corners=3pt, fill=black!4,
      inner sep=4pt, font=\scriptsize, anchor=north] at (4.2,-1.4)
     {$q = \dfrac{\text{amplitude du signal d'entrée}}{2^{n}}$};
\end{tikzpicture}
\end{document}
